% Seminar 3 - Page 4. Fragment for
% % Seminar 3 - Page 4. Fragment for
% % Seminar 3 - Page 4. Fragment for
% % Seminar 3 - Page 4. Fragment for
% \input{seminars/pages/seminar3-pagina4-en}; no preamble or page break.
% EDITORIAL NOTES:
% - The first three lines continue Exercise 2 from Page 3; the proof is
%   not presented as a new exercise here.
% - The handwritten cardinality of im(f) intersect [n,n+1) reads <= 1.
% - The equivalence relation and representative-system symbols are
%   re-set for legibility; their mathematical content is preserved.
% - The manuscript gives the properties of a representative system
%   compactly, with S written beneath (x_i)_{i in I}.
% - The diagram of the quotient's universal property is redrawn as a
%   standard commutative triangle, preserving the named arrows.
% - The universal-property theorem is stated, not proved, on this page.

\subsection*{Seminar 3 - Page 4}
\begingroup
\setlength{\abovedisplayskip}{5pt}
\setlength{\belowdisplayskip}{5pt}
\setlength{\abovedisplayshortskip}{3pt}
\setlength{\belowdisplayshortskip}{3pt}

\noindent\textbf{Exercise 2 (continued).} From the previous page,
\[
 \operatorname{im}(f)=\bigcup_{n\in\mathbb Z}
 \bigl(\operatorname{im}(f)\cap[n,n+1)\bigr).
\]
Since $|f(x)-f(y)|>1$ for $x\neq y$, we have
$|\operatorname{im}(f)\cap[n,n+1)|\leq1$.
Hence $\operatorname{im}(f)$ is finite or countable: a contradiction.

\medskip
\noindent\underline{Equivalence relations}

\noindent A relation $R\subseteq M\times M$ is called an
\emph{equivalence relation} if it satisfies:
\[
\begin{array}{ll}
 \text{(R) Reflexivity:}&(x,x)\in R,\\[2pt]
 \text{(S) Symmetry:}&(x,y)\in R\ \Longrightarrow\ (y,x)\in R,\\[2pt]
 \text{(T) Transitivity:}&
  \left.\begin{array}{l}(x,y)\in R\\(y,z)\in R\end{array}\right\}
  \Longrightarrow (x,z)\in R.
\end{array}
\]
Write $x\sim y$ when $(x,y)\in R$. The equivalence class of
$x\in M$ is
\[
 \widehat{x}=\{y\in M\mid y\sim x\}.
\]
The quotient set and canonical projection are
\[
 M/{\sim}:=\{\widehat{x}\mid x\in M\},
 \qquad
 \pi:M\longrightarrow M/{\sim},
 \quad \pi(x)=\widehat{x}.
\]
\noindent 
\noindent A \emph{system of representatives} is a family
$S=(x_i)_{i\in I}$ such that
\[
 \begin{cases}
   x_i\not\sim x_j,&i\neq j,\\[2pt]
   \forall x\in M\ \exists i\in I:\ x\sim x_i.
 \end{cases}
\]
For the corresponding subset $S\subseteq M$, the restriction
$\pi|_S:S\to M/{\sim}$ is bijective.

\medskip
\noindent\underline{Universal property}

\noindent\textbf{Observation.} Given $f:M\to X$, define
\[
 R_f:=\{(x,y)\in M\times M\mid f(x)=f(y)\}.
\]
Then $x\sim y\iff f(x)=f(y)$ defines the equivalence relation
induced by $f$ on $M$.

\noindent\fbox{T}\quad \textbf{Universal property of the quotient set.}
Let $\sim$ be an equivalence relation on $M$ and let
$\pi:M\to M/{\sim}$ be the canonical projection.
For any map $f:M\to X$ satisfying $\sim\,\subseteq R_f$,
there exists a unique map $\overline f:M/{\sim}\to X$ such that
\[
 \overline f\circ\pi=f.
\]
\[
\begin{tikzcd}[column sep=large,row sep=large]
 M\arrow[r,"\pi"]\arrow[dr,"f"']
   & M/{\sim}\arrow[d,dashed,"{\overline f}"]\\
 & X
\end{tikzcd}
\]
% The manuscript stops after the statement of the universal property.

\endgroup
; no preamble or page break.
% EDITORIAL NOTES:
% - The first three lines continue Exercise 2 from Page 3; the proof is
%   not presented as a new exercise here.
% - The handwritten cardinality of im(f) intersect [n,n+1) reads <= 1.
% - The equivalence relation and representative-system symbols are
%   re-set for legibility; their mathematical content is preserved.
% - The manuscript gives the properties of a representative system
%   compactly, with S written beneath (x_i)_{i in I}.
% - The diagram of the quotient's universal property is redrawn as a
%   standard commutative triangle, preserving the named arrows.
% - The universal-property theorem is stated, not proved, on this page.

\subsection*{Seminar 3 - Page 4}
\begingroup
\setlength{\abovedisplayskip}{5pt}
\setlength{\belowdisplayskip}{5pt}
\setlength{\abovedisplayshortskip}{3pt}
\setlength{\belowdisplayshortskip}{3pt}

\noindent\textbf{Exercise 2 (continued).} From the previous page,
\[
 \operatorname{im}(f)=\bigcup_{n\in\mathbb Z}
 \bigl(\operatorname{im}(f)\cap[n,n+1)\bigr).
\]
Since $|f(x)-f(y)|>1$ for $x\neq y$, we have
$|\operatorname{im}(f)\cap[n,n+1)|\leq1$.
Hence $\operatorname{im}(f)$ is finite or countable: a contradiction.

\medskip
\noindent\underline{Equivalence relations}

\noindent A relation $R\subseteq M\times M$ is called an
\emph{equivalence relation} if it satisfies:
\[
\begin{array}{ll}
 \text{(R) Reflexivity:}&(x,x)\in R,\\[2pt]
 \text{(S) Symmetry:}&(x,y)\in R\ \Longrightarrow\ (y,x)\in R,\\[2pt]
 \text{(T) Transitivity:}&
  \left.\begin{array}{l}(x,y)\in R\\(y,z)\in R\end{array}\right\}
  \Longrightarrow (x,z)\in R.
\end{array}
\]
Write $x\sim y$ when $(x,y)\in R$. The equivalence class of
$x\in M$ is
\[
 \widehat{x}=\{y\in M\mid y\sim x\}.
\]
The quotient set and canonical projection are
\[
 M/{\sim}:=\{\widehat{x}\mid x\in M\},
 \qquad
 \pi:M\longrightarrow M/{\sim},
 \quad \pi(x)=\widehat{x}.
\]
\noindent 
\noindent A \emph{system of representatives} is a family
$S=(x_i)_{i\in I}$ such that
\[
 \begin{cases}
   x_i\not\sim x_j,&i\neq j,\\[2pt]
   \forall x\in M\ \exists i\in I:\ x\sim x_i.
 \end{cases}
\]
For the corresponding subset $S\subseteq M$, the restriction
$\pi|_S:S\to M/{\sim}$ is bijective.

\medskip
\noindent\underline{Universal property}

\noindent\textbf{Observation.} Given $f:M\to X$, define
\[
 R_f:=\{(x,y)\in M\times M\mid f(x)=f(y)\}.
\]
Then $x\sim y\iff f(x)=f(y)$ defines the equivalence relation
induced by $f$ on $M$.

\noindent\fbox{T}\quad \textbf{Universal property of the quotient set.}
Let $\sim$ be an equivalence relation on $M$ and let
$\pi:M\to M/{\sim}$ be the canonical projection.
For any map $f:M\to X$ satisfying $\sim\,\subseteq R_f$,
there exists a unique map $\overline f:M/{\sim}\to X$ such that
\[
 \overline f\circ\pi=f.
\]
\[
\begin{tikzcd}[column sep=large,row sep=large]
 M\arrow[r,"\pi"]\arrow[dr,"f"']
   & M/{\sim}\arrow[d,dashed,"{\overline f}"]\\
 & X
\end{tikzcd}
\]
% The manuscript stops after the statement of the universal property.

\endgroup
; no preamble or page break.
% EDITORIAL NOTES:
% - The first three lines continue Exercise 2 from Page 3; the proof is
%   not presented as a new exercise here.
% - The handwritten cardinality of im(f) intersect [n,n+1) reads <= 1.
% - The equivalence relation and representative-system symbols are
%   re-set for legibility; their mathematical content is preserved.
% - The manuscript gives the properties of a representative system
%   compactly, with S written beneath (x_i)_{i in I}.
% - The diagram of the quotient's universal property is redrawn as a
%   standard commutative triangle, preserving the named arrows.
% - The universal-property theorem is stated, not proved, on this page.

\subsection*{Seminar 3 - Page 4}
\begingroup
\setlength{\abovedisplayskip}{5pt}
\setlength{\belowdisplayskip}{5pt}
\setlength{\abovedisplayshortskip}{3pt}
\setlength{\belowdisplayshortskip}{3pt}

\noindent\textbf{Exercise 2 (continued).} From the previous page,
\[
 \operatorname{im}(f)=\bigcup_{n\in\mathbb Z}
 \bigl(\operatorname{im}(f)\cap[n,n+1)\bigr).
\]
Since $|f(x)-f(y)|>1$ for $x\neq y$, we have
$|\operatorname{im}(f)\cap[n,n+1)|\leq1$.
Hence $\operatorname{im}(f)$ is finite or countable: a contradiction.

\medskip
\noindent\underline{Equivalence relations}

\noindent A relation $R\subseteq M\times M$ is called an
\emph{equivalence relation} if it satisfies:
\[
\begin{array}{ll}
 \text{(R) Reflexivity:}&(x,x)\in R,\\[2pt]
 \text{(S) Symmetry:}&(x,y)\in R\ \Longrightarrow\ (y,x)\in R,\\[2pt]
 \text{(T) Transitivity:}&
  \left.\begin{array}{l}(x,y)\in R\\(y,z)\in R\end{array}\right\}
  \Longrightarrow (x,z)\in R.
\end{array}
\]
Write $x\sim y$ when $(x,y)\in R$. The equivalence class of
$x\in M$ is
\[
 \widehat{x}=\{y\in M\mid y\sim x\}.
\]
The quotient set and canonical projection are
\[
 M/{\sim}:=\{\widehat{x}\mid x\in M\},
 \qquad
 \pi:M\longrightarrow M/{\sim},
 \quad \pi(x)=\widehat{x}.
\]
\noindent 
\noindent A \emph{system of representatives} is a family
$S=(x_i)_{i\in I}$ such that
\[
 \begin{cases}
   x_i\not\sim x_j,&i\neq j,\\[2pt]
   \forall x\in M\ \exists i\in I:\ x\sim x_i.
 \end{cases}
\]
For the corresponding subset $S\subseteq M$, the restriction
$\pi|_S:S\to M/{\sim}$ is bijective.

\medskip
\noindent\underline{Universal property}

\noindent\textbf{Observation.} Given $f:M\to X$, define
\[
 R_f:=\{(x,y)\in M\times M\mid f(x)=f(y)\}.
\]
Then $x\sim y\iff f(x)=f(y)$ defines the equivalence relation
induced by $f$ on $M$.

\noindent\fbox{T}\quad \textbf{Universal property of the quotient set.}
Let $\sim$ be an equivalence relation on $M$ and let
$\pi:M\to M/{\sim}$ be the canonical projection.
For any map $f:M\to X$ satisfying $\sim\,\subseteq R_f$,
there exists a unique map $\overline f:M/{\sim}\to X$ such that
\[
 \overline f\circ\pi=f.
\]
\[
\begin{tikzcd}[column sep=large,row sep=large]
 M\arrow[r,"\pi"]\arrow[dr,"f"']
   & M/{\sim}\arrow[d,dashed,"{\overline f}"]\\
 & X
\end{tikzcd}
\]
% The manuscript stops after the statement of the universal property.

\endgroup
; no preamble or page break.
% EDITORIAL NOTES:
% - The first three lines continue Exercise 2 from Page 3; the proof is
%   not presented as a new exercise here.
% - The handwritten cardinality of im(f) intersect [n,n+1) reads <= 1.
% - The equivalence relation and representative-system symbols are
%   re-set for legibility; their mathematical content is preserved.
% - The manuscript gives the properties of a representative system
%   compactly, with S written beneath (x_i)_{i in I}.
% - The diagram of the quotient's universal property is redrawn as a
%   standard commutative triangle, preserving the named arrows.
% - The universal-property theorem is stated, not proved, on this page.

\subsection*{Seminar 3 - Page 4}
\begingroup
\setlength{\abovedisplayskip}{5pt}
\setlength{\belowdisplayskip}{5pt}
\setlength{\abovedisplayshortskip}{3pt}
\setlength{\belowdisplayshortskip}{3pt}

\noindent\textbf{Exercise 2 (continued).} From the previous page,
\[
 \operatorname{im}(f)=\bigcup_{n\in\mathbb Z}
 \bigl(\operatorname{im}(f)\cap[n,n+1)\bigr).
\]
Since $|f(x)-f(y)|>1$ for $x\neq y$, we have
$|\operatorname{im}(f)\cap[n,n+1)|\leq1$.
Hence $\operatorname{im}(f)$ is finite or countable: a contradiction.

\medskip
\noindent\underline{Equivalence relations}

\noindent A relation $R\subseteq M\times M$ is called an
\emph{equivalence relation} if it satisfies:
\[
\begin{array}{ll}
 \text{(R) Reflexivity:}&(x,x)\in R,\\[2pt]
 \text{(S) Symmetry:}&(x,y)\in R\ \Longrightarrow\ (y,x)\in R,\\[2pt]
 \text{(T) Transitivity:}&
  \left.\begin{array}{l}(x,y)\in R\\(y,z)\in R\end{array}\right\}
  \Longrightarrow (x,z)\in R.
\end{array}
\]
Write $x\sim y$ when $(x,y)\in R$. The equivalence class of
$x\in M$ is
\[
 \widehat{x}=\{y\in M\mid y\sim x\}.
\]
The quotient set and canonical projection are
\[
 M/{\sim}:=\{\widehat{x}\mid x\in M\},
 \qquad
 \pi:M\longrightarrow M/{\sim},
 \quad \pi(x)=\widehat{x}.
\]
\noindent 
\noindent A \emph{system of representatives} is a family
$S=(x_i)_{i\in I}$ such that
\[
 \begin{cases}
   x_i\not\sim x_j,&i\neq j,\\[2pt]
   \forall x\in M\ \exists i\in I:\ x\sim x_i.
 \end{cases}
\]
For the corresponding subset $S\subseteq M$, the restriction
$\pi|_S:S\to M/{\sim}$ is bijective.

\medskip
\noindent\underline{Universal property}

\noindent\textbf{Observation.} Given $f:M\to X$, define
\[
 R_f:=\{(x,y)\in M\times M\mid f(x)=f(y)\}.
\]
Then $x\sim y\iff f(x)=f(y)$ defines the equivalence relation
induced by $f$ on $M$.

\noindent\fbox{T}\quad \textbf{Universal property of the quotient set.}
Let $\sim$ be an equivalence relation on $M$ and let
$\pi:M\to M/{\sim}$ be the canonical projection.
For any map $f:M\to X$ satisfying $\sim\,\subseteq R_f$,
there exists a unique map $\overline f:M/{\sim}\to X$ such that
\[
 \overline f\circ\pi=f.
\]
\[
\begin{tikzcd}[column sep=large,row sep=large]
 M\arrow[r,"\pi"]\arrow[dr,"f"']
   & M/{\sim}\arrow[d,dashed,"{\overline f}"]\\
 & X
\end{tikzcd}
\]
% The manuscript stops after the statement of the universal property.

\endgroup
